%---------------------------------------------------------------------------------------
%	Honors \& Awards
%--------------------------------------------------------------------------------------
\cvsection{Research Projects}


\cvevent
{2020 - Present}
{Conversational News}
{Martin-Luther-University, Bauhaus-Universität Weimar} % 
\cvdetail{This project includes the development of a prototype that deploy text news articles to smart assistants. Writing augmented tools will be provided by semi-automatic ways to enhance the experience of delivering news in audio channels.}
\cvdetail{The project is funded by Google within the Digital News Innovation Fund.}

\cvevent
{2019 - Present}
{Argumentation Knowledge Graph}
{Bauhaus-Universität Weimar (Webis.de)} %
\cvdetail {The project aims for constructing an argumentation knowledge graph that is adequate for supporting several artificial intelligent down-stream applications such as text generation, decision making, and argument search.}
\cvdetail{The project is partially funded by IBM through the Ph.D. fellowship award.}

\cvevent
{2015 - 2018}
{Args.me Search Engine}
{Bauhaus-Universität Weimar (Webis.de)} %
\cvdetail {This project targets the development of the first argument search engine on the web. The project addresses several challenges of mining arguments and their relations from natural language text on the web. We develop algorithms that (1) robustly tackle various forms of web argumentation, (2) efficiently leverage the scale of the web, and (3) effectively complement argument mining with an argumentation analysis to assess important quality dimensions.}
\cvdetail{The project is funded by the German Research Foundation (DFG).}
\cvdetail{Advisor: Prof. Dr. Benno Stein. Prof. Dr. Matthias Hagen, Jun. Prof. Dr. Henning Wachsmuth}

\cvevent
{2015 - Present}
{Argumentation Strategies}% Job title
{Bauhaus-Universität Weimar (Webis.de)} %
\cvdetail {This research project is centered around the strategies that authors of monological argumentative texts, such as news editorials and scientific articles, employ to persuade their readers. It seeks to (1)~establish computational models of such argumentation strategies, (2)~analyze what strategies are common and effective within and across argumentative text genres, and (3)~synthesize texts with effective strategies for given sets of arguments.}
\cvdetail{The project is funded by the German Research Foundation (DFG).}
\cvdetail{Advisor: Prof.Dr. Benno Stein. Jun.-Prof. Dr. Henning Wachsmuth.}

\cvevent
{11.2018 - 04.2019} % Date(s)} 
{Counter-argument Generation} 
{IBM Research, Ireland Research Lab}
\cvdetail {This project concerns developing new technologies for tackling the following task:given a target argument, we aim to recommend its counterargument(s) from a set of mined arguments related to the same topic.}
\cvdetail{Advisor: Dr. Lea Deleris. Dr. Yufang Hou}

\cvevent
{2011-2013} 
{Sentiment Topic Analysis} 
{Institute for NLP, University of Stuttgart} 
\cvdetail {This project studies some areas in sentiment analysis that haven’t been studied extensively before like perspectives’ classification, bias detection, and the relation between social networks and their users with the public opinion of different sub-communities.}
\cvdetail{Advisor: Prof. Dr. Hinrich Schutze.}

\cvevent
{2011-2013} 
{eIdentity} 
{Institute for NLP, University of Stuttgart} 
\cvdetail {This project studies multiple collective identities in international debates about war and peace since the end of the Cold War. Many language technology tools
and methods are used to analyze multilingual text concepts in the social sciences.}
%My Role: I work on mining Wikipedia tables and extracting of temporal facts about politicians.\\
\cvdetail{Advisor: Prof. Dr. Cathleen Kantner.}
